\section{Introduction}\label{sec:intro}

Consider a persistent AI companion. A user raises a digital character over
hundreds of turns---practitioners sometimes call this ``cyber fish-keeping,''
slang rather than a research construct---and one day the character forgets its
own name or catchphrase; close the session, and it is gone. The motivation for
this work began as a simple wish: not wanting a virtual companion to die.

An interactive agent is not a single model call: before every call the host
assembles a prompt from memories, persona, goals, affect, and safety policy. In
short sessions this assembly is invisible; in long sessions it becomes the
dominant variable, because injected material grows while the budget does not.

Four failure modes follow, and they are the object of design throughout the
paper. \textbf{F1, memory overreach}: a stored memory is presented as a
statement of fact rather than as a reference to something the user once said.
\textbf{F2, refusal softening}: stylistic or affective signals are injected
into a turn that should be a refusal, and the refusal is weakened.
\textbf{F3, persona drift}: persona constraints stop influencing behaviour as
the context fills. \textbf{F4, budget capture}: truncation removes a
high-priority block while lower-priority stylistic blocks survive. Each is
observable at a specific point of the assembly pipeline, which makes them
tractable. We claim a control surface that is predictable, explainable and
rollbackable, and state its limits in Section~7.

\paragraph{Why structural, not only behavioural.} Alignment training shapes model
behaviour probabilistically: it lowers the rate of unsafe responses but does not
constrain what the host assembles around the model, nor certify after the fact
that a turn respected the policy. A structural layer complements it: decisions
are made in code at the assembly boundary before the model call, each decision is
recorded and replayable, and the one component whose failure is a safety failure
is fail-closed. The goal is not a better-behaved model but a prompt-assembly
surface whose behaviour can be constrained, observed and rolled back.

\paragraph{Positioning.} \system{} is not an attempt to make an agent more
human. Four senses of ``human-like'' are commonly conflated: capability
(remembering, forgetting, tracking goals), imperfection (misremembering, being
influenced by stale information), expression (tone, empathy, level of detail),
and self-claim (asserting feelings, needs, or subjective experience). We target
the first three and reject the fourth: it conflicts with the honesty property
the system should provide. The design goal is narrower: make
agent behaviour understandable, predictable, and trustworthy to the human in the
loop.

\paragraph{Five invariants.} Every component is constrained by five
invariants. \textbf{I1}: memory is a reference and never replaces the user's
current statement. \textbf{I2}: no full-history injection; context is a budget,
not a dump. \textbf{I3}: affect influences style only, never factual judgement.
\textbf{I4}: every inference is verbatim traceable, deniable, and rollbackable.
\textbf{I5}: plugins write only their own storage, and the memory plugin is the
single writer. A violation is a defect, not a trade-off.

\paragraph{Naming the control surface.} The mechanisms that police prompt
assembly---integral feedback, preemptive monitoring, dimensionality reduction,
and the pre-flight gate---are collectively what we call the
\emph{cyber-cognitive envelope}. The name is a framing device, not a new
mechanism: it denotes no separate code entity, each part is an ordinary module,
and no claim below depends on it.

\paragraph{Contributions.} Our contributions are structural and checkable; we
order them by how directly each one constrains the assembly boundary.

\begin{itemize}
\item \textbf{C1, a machine-checkable contract} (directly rerunnable). The
specification is not a document but assertions that fail loudly when a contract
breaks: 892 Python assertions across 14 test files, plus 396 plugin-side
assertions (188 host integration, 148 control-plane unit, 45 stress). Parameter
consistency between rules, engine, and plugin copies is itself machine-checked
(66 key comparisons).
\item \textbf{C2, one fail-closed locus} (pinned by regression tests). Exactly
one component is fail-closed---the STPA audit; everything else degrades open and
records why. Eight unsafe control actions are screened after assembly and before
the model call, and integral saturation is recorded as a downgrade rather than a
block. Mutation testing on the gate and memory-write paths kills
$42/42=100\%$ of non-equivalent mutants, with two equivalent mutants reported as
such.
\item \textbf{C3, structural defence at the assembly boundary} (measured, with a
stated boundary). Five domains are composed under a frozen order and eight
mutual-exclusion rules; what survives composition is screened by the gate. The
envelope therefore structurally constrains what reaches the model and fails
closed on violation, rather than asking the model to comply. Where its lexical
coverage ends is measured rather than assumed (Section~7).
\item \textbf{C4, an evaluation discipline and a documented port} (artefacts
exist and are machine-checked). Every parameter carries a provenance tag; the
instrument (metric definitions, replay set, guardrail thresholds) is frozen and
read-only for the learning loop; changes to it are rejected by an auditor. The
same control semantics exist twice, in a zero-dependency Python reference
implementation and a JavaScript host plugin, and the five known divergences are
listed rather than hidden.
\end{itemize}

\paragraph{Scope of the evidence.} The paper reports two kinds of evidence.
Contract feasibility: the specification executes and violation is observable.
Mechanism reachability: in the tested configuration each mechanism fires in a
real host session, with sample sizes stated. It also reports one negative result
that shaped the design: replacing exponential decay with a pure ACT-R base-level
activation drops \texttt{recall@3} on the frozen replay set from 1.000 to 0.556,
which is why the implemented model is a dual-channel maximum rather than a
replacement. Effect sizes on real traffic and the calibration of design-default
parameters are outside this paper's scope; Section~7 states what the current
evidence supports.

Section~\ref{sec:related} reviews related work, Section~\ref{sec:problem}
formalizes the problem, Sections~4--5 present mechanisms and evidence, and
Sections~6--8 discuss costs, limits, and conclusions.
